% Generated by publication/build_sections.py; do not edit by hand.
\paragraph{Informal verdict.}
\textbf{A --- Complete and apparently correct.} Confidence: \textbf{high (0.96)}.

\paragraph{Lean coverage.}
Lean proves at p=4, with C=2*sqrt\(2\)>=1, the actual square-function estimate uniformly for every N (and hence every positive N), every Fin N family of nondegenerate closed subintervals of [0,1], and every corresponding family of real L4 functions; Q899.exists\_exponent\_gt\_two exports the printed existential statement for every positive N. The larger constant than the informal 2\textasciicircum{}(3/4) is canonically sufficient.

\subsection{Audited informal answer}
Q899 --- Complete solution
\begin{enumerate}
\item Verified statement
\end{enumerate}
For every nondegenerate interval \(I=[a,b]\subset[0,1]\), write \(|I|=b-a\) and define the rank-one projection

\[
P_I f
=
\left(\frac1{|I|}\int_I f(y)\,dy\right)\mathbf 1_I .
\]

The problem asks whether there exist \(p>2\) and \(C(p)\ge 1\) such that, for every \(N\ge1\), every choice of intervals \(I_1,\dots,I_N\subset[0,1]\), and every \(f_1,\dots,f_N\in L^p(0,1)\),

\[
\left\|
\left(\sum_{j=1}^N |P_{I_j}f_j|^2\right)^{1/2}
\right\|_{L^p(0,1)}
\le
C(p)
\left\|
\left(\sum_{j=1}^N |f_j|^2\right)^{1/2}
\right\|_{L^p(0,1)}.
\tag{1}
\]

The PDF describes this as the \(R\)-boundedness of the family \(\{P_I\}\). Luc Pommeret
Two minor formal clarifications are implicit in the statement:


One must require \(a<b\), since the displayed definition contains \((b-a)^{-1}\).


Since \(f\in L^p\) is an equivalence class, the defining equality for \(P_I f\) is understood almost everywhere, rather than literally for every \(x\).


There is no other ambiguity.

\begin{enumerate}
\item Final theorem
\end{enumerate}
Let \(\mathfrak R_p\) denote the least constant in \(1\), uniformly over all finite families of intervals and functions.
Theorem
The family of interval projections satisfies

\[
\mathfrak R_p<\infty
\quad\Longleftrightarrow\quad
1<p<\infty.
\]

More precisely:


For every \(p\ge2\),

\[
\boxed{\mathfrak R_p\le 2^{-1/p}\sqrt p.}
\tag{2}
\]



At \(p=2\),

\[
\boxed{\mathfrak R_2=1.}
\]



If \(1<p<2\) and \(p'=p/(p-1)\), then

\[
\mathfrak R_p=\mathfrak R_{p'}
\le 2^{-1/p'}\sqrt{p'}.
\tag{3}
\]



The inequality fails uniformly for \(p=1\) and \(p=\infty\).


Consequently, the answer to Q899 is yes for every \(p>2\). For example, one may take

\[
\boxed{p=4,\qquad C(4)=2^{3/4}.}
\]

The constant displayed here is not asserted to be optimal.

\begin{enumerate}
\item A scalar maximal-function estimate
\end{enumerate}
For \(h\in L^1_{\mathrm{loc}}(\mathbb R)\), define the one-sided maximal functions

\[
M^+h(x)=\sup_{t>0}\frac1t\int_x^{x+t}|h(y)|\,dy,
\]


\[
M^-h(x)=\sup_{t>0}\frac1t\int_{x-t}^{x}|h(y)|\,dy.
\]

We first prove the estimate needed below without invoking the vector-valued maximal theorem.
Lemma 1
For \(1<q<\infty\), with \(q'=q/(q-1)\),

\[
\|M^+h\|_{L^q(\mathbb R)}
\le q'\|h\|_{L^q(\mathbb R)}
\tag{4}
\]

and similarly

\[
\|M^-h\|_{L^q(\mathbb R)}
\le q'\|h\|_{L^q(\mathbb R)}.
\tag{5}
\]

Proof
It suffices to consider \(h\ge0\). We first suppose that \(h\) is bounded and compactly supported.
Fix \(\lambda>0\), and define

\[
F_\lambda(x)=\int_{-\infty}^x h(t)\,dt-\lambda x.
\]

Then

\[
M^+h(x)>\lambda
\]

if and only if there exists \(y>x\) such that

\[
\frac1{y-x}\int_x^y h(t)\,dt>\lambda,
\]

or equivalently,

\[
F_\lambda(y)>F_\lambda(x).
\]

Set

\[
E_\lambda=\{x\in\mathbb R:M^+h(x)>\lambda\}.
\]

The set \(E_\lambda\) is open and bounded, hence is a countable disjoint union of bounded open intervals,

\[
E_\lambda=\bigcup_k(a_k,b_k).
\]

We claim that

\[
F_\lambda(a_k)=F_\lambda(b_k)
\tag{6}
\]

for every component.
Indeed, since \(a_k\notin E_\lambda\), no later point has strictly larger \(F_\lambda\)-value, and therefore

\[
F_\lambda(b_k)\le F_\lambda(a_k).
\tag{7}
\]

Conversely, take \(x\in(a_k,b_k)\). Since \(F_\lambda(y)\to-\infty\) as \(y\to+\infty\), \(F_\lambda\) attains its maximum on \([x,\infty)\). Choose the largest point \(y_x\) where this maximum is attained. Then \(y_x\notin E_\lambda\); hence, because \((a_k,b_k)\) is a connected component containing \(x\), one has \(y_x\ge b_k\). Since \(b_k\notin E_\lambda\), no point after \(b_k\) has value larger than \(F_\lambda(b_k)\). Consequently,

\[
F_\lambda(y_x)=F_\lambda(b_k).
\]

Because \(x\in E_\lambda\),

\[
F_\lambda(x)<F_\lambda(y_x)=F_\lambda(b_k).
\]

Letting \(x\downarrow a_k\) and using continuity of \(F_\lambda\) gives

\[
F_\lambda(a_k)\le F_\lambda(b_k).
\]

Together with \(7\), this proves \(6\). Therefore,

\[
\int_{a_k}^{b_k}h(t)\,dt
=
\lambda(b_k-a_k).
\]

Summing over the components gives the localized weak-type identity

\[
\lambda |E_\lambda|
=
\int_{E_\lambda}h(x)\,dx.
\tag{8}
\]

Since \(M^+h\le\|h\|_\infty\), the layer-cake formula and Tonelli's theorem yield

\[
\begin{aligned}
\|M^+h\|_q^q
&=
q\int_0^\infty
\lambda^{q-1}
|\{M^+h>\lambda\}|\,d\lambda\\
&=
q\int_0^\infty
\lambda^{q-2}
\int_{\{M^+h>\lambda\}}h(x)\,dx\,d\lambda\\
&=
q\int_{\mathbb R}
h(x)
\int_0^{M^+h(x)}
\lambda^{q-2}\,d\lambda\,dx\\
&=
\frac{q}{q-1}
\int_{\mathbb R}
h(x)\bigl(M^+h(x)\bigr)^{q-1}\,dx.
\end{aligned}
\]

By Hölder's inequality,

\[
\|M^+h\|_q^q
\le
q'\|h\|_q
\|M^+h\|_q^{q-1}.
\]

If \(\|M^+h\|_q\ne0\), division gives \(4\).
For a general nonnegative \(h\in L^q(\mathbb R)\), take

\[
h_n=\min(h,n)\mathbf 1_{[-n,n]}.
\]

Then \(h_n\uparrow h\) and \(M^+h_n\uparrow M^+h\). Applying the preceding estimate to \(h_n\) and using monotone convergence proves \(4\) in general.
Estimate \(5\) follows by applying \(4\) to \(x\mapsto h(-x)\). \(\square\)
Now define the uncentered maximal function on \([0,1]\) by

\[
\mathcal Mh(x)
=
\sup_{\substack{I\subset[0,1]\\x\in I}}
\frac1{|I|}\int_I |h(y)|\,dy.
\]

Extend \(h\) by zero outside \([0,1]\). If \(I=[a,b]\ni x\), then its average is a convex combination of the left and right averages:

\[
\frac1{b-a}\int_a^b |h|
=
\frac{x-a}{b-a}
\frac1{x-a}\int_a^x |h|
+
\frac{b-x}{b-a}
\frac1{b-x}\int_x^b |h|.
\]

Thus, almost everywhere,

\[
\mathcal Mh(x)\le\max\{M^-h(x),M^+h(x)\}.
\]

Consequently, for \(1<q<\infty\),

\[
\begin{aligned}
\|\mathcal Mh\|_q^q
&\le
\|M^-h\|_q^q+\|M^+h\|_q^q\\
&\le
2(q')^q\|h\|_q^q,
\end{aligned}
\]

and hence

\[
\boxed{
\|\mathcal Mh\|_{L^q(0,1)}
\le
2^{1/q}q'\|h\|_{L^q(0,1)}.
}
\tag{9}
\]


\begin{enumerate}
\item The square-function inequality for \(p>2\)
\end{enumerate}
Fix \(p>2\), and put

\[
r=\frac p2,
\qquad
q=r'=\frac{p}{p-2}.
\]

Thus \(r,q\in(1,\infty)\) and

\[
q'=\frac q{q-1}=\frac p2.
\]

Let \(I_j=[a_j,b_j]\) and set

\[
m_j=\frac1{|I_j|}\int_{I_j}f_j(y)\,dy.
\]

Then

\[
P_{I_j}f_j=m_j\mathbf 1_{I_j}.
\]

Define

\[
F(x)=
\left(\sum_{j=1}^N|f_j(x)|^2\right)^{1/2},
\qquad
G(x)=
\left(\sum_{j=1}^N|P_{I_j}f_j(x)|^2\right)^{1/2}.
\]

Since \(p>2\) and \([0,1]\) has finite measure, each \(f_j\in L^2(0,1)\), so all the following integrals are finite.
Because \(G^2\ge0\), duality in \(L^r\) gives

\[
\|G\|_p^2
=
\|G^2\|_r
=
\sup_{\substack{h\ge0\\ \|h\|_q=1}}
\int_0^1 G(x)^2h(x)\,dx.
\tag{10}
\]

Fix such an \(h\). By Cauchy-Schwarz,

\[
|m_j|^2
=
\left\lvert{}
\frac1{|I_j|}\int_{I_j}f_j
\right\rvert{}^2
\le
\frac1{|I_j|}\int_{I_j}|f_j(y)|^2\,dy.
\tag{11}
\]

Therefore,

\[
\begin{aligned}
\int_0^1G^2h
&=
\sum_{j=1}^N
|m_j|^2\int_{I_j}h(x)\,dx\\
&\le
\sum_{j=1}^N
\left(
\frac1{|I_j|}
\int_{I_j}|f_j(y)|^2\,dy
\right)
\int_{I_j}h(x)\,dx\\
&=
\sum_{j=1}^N
\int_{I_j}
|f_j(y)|^2
\left(
\frac1{|I_j|}\int_{I_j}h(x)\,dx
\right)dy.
\end{aligned}
\]

For \(y\in I_j\),

\[
\frac1{|I_j|}\int_{I_j}h
\le \mathcal Mh(y).
\]

Hence

\[
\int_0^1G^2h
\le
\int_0^1
\left(\sum_{j=1}^N|f_j(y)|^2\right)
\mathcal Mh(y)\,dy
=
\int_0^1F(y)^2\mathcal Mh(y)\,dy.
\]

Hölder's inequality with exponents \(r=p/2\) and \(q=p/(p-2)\), followed by \(9\), yields

\[
\begin{aligned}
\int_0^1G^2h
&\le
\|F^2\|_r\|\mathcal Mh\|_q\\
&\le
2^{1/q}q'\|F^2\|_r\|h\|_q\\
&=
2^{1/q}q'\|F\|_p^2.
\end{aligned}
\]

Since \(\|h\|_q=1\), taking the supremum in \(10\) gives

\[
\|G\|_p^2
\le
2^{1/q}q'\|F\|_p^2.
\]

Now

\[
\frac1q=\frac{p-2}{p},
\qquad
q'=\frac p2,
\]

so

\[
2^{1/q}q'
=
2^{(p-2)/p}\frac p2
=
p\,2^{-2/p}.
\]

Taking square roots,

\[
\boxed{
\|G\|_p
\le
2^{-1/p}\sqrt p\,\|F\|_p.
}
\tag{12}
\]

This proves the required inequality for every \(p>2\).
For instance, when \(p=4\),

\[
2^{-1/4}\sqrt4=2^{3/4},
\]

so \(p=4\) and \(C(4)=2^{3/4}\) give an explicit answer to the original existence question.

\begin{enumerate}
\item The Hilbert-space boundary \(p=2\)
\end{enumerate}
For \(p=2\), directly from \(11\),

\[
\begin{aligned}
\left\|
\left(\sum_{j=1}^N|P_{I_j}f_j|^2\right)^{1/2}
\right\|_2^2
&=
\sum_{j=1}^N
\|P_{I_j}f_j\|_2^2\\
&=
\sum_{j=1}^N
|m_j|^2|I_j|\\
&\le
\sum_{j=1}^N
\int_{I_j}|f_j(y)|^2\,dy\\
&\le
\int_0^1\sum_{j=1}^N|f_j(y)|^2\,dy.
\end{aligned}
\]

Thus \(\mathfrak R_2\le1\). Equality is necessary already for \(N=1\) and \(f=\mathbf1_I\), because \(P_I f=f\). Therefore

\[
\boxed{\mathfrak R_2=1.}
\]

Notice that the upper bound \(2^{-1/p}\sqrt p\) tends to \(1\) as \(p\downarrow2\), consistently with this exact value.

\begin{enumerate}
\item Extension to \(1<p<2\) by duality
\end{enumerate}
For fixed intervals \(I_1,\dots,I_N\), define

\[
T_{\mathcal I}:
L^p\bigl((0,1);\ell_2^N\bigr)
\longrightarrow
L^p\bigl((0,1);\ell_2^N\bigr)
\]

by

\[
T_{\mathcal I}(f_1,\dots,f_N)
=
(P_{I_1}f_1,\dots,P_{I_N}f_N).
\]

The inequality under consideration is exactly the uniform boundedness of these operators.
Each \(P_I\) is self-adjoint with respect to the integral pairing. Indeed,

\[
\begin{aligned}
\int_0^1 P_I f(x)\,\overline{g(x)}\,dx
&=
\frac1{|I|}
\left(\int_I f\right)
\left(\int_I\overline g\right)\\
&=
\int_0^1
f(x)\,\overline{P_Ig(x)}\,dx.
\end{aligned}
\]

Therefore, under the duality

\[
L^p(\ell_2^N)^*=L^{p'}(\ell_2^N),
\]

one has

\[
T_{\mathcal I}^*=T_{\mathcal I}.
\]

Consequently,

\[
\|T_{\mathcal I}\|_{L^p(\ell_2^N)\to L^p(\ell_2^N)}
=
\|T_{\mathcal I}\|_{L^{p'}(\ell_2^N)\to L^{p'}(\ell_2^N)}.
\]

Taking the supremum over \(N\) and over all interval families gives

\[
\boxed{\mathfrak R_p=\mathfrak R_{p'}.}
\tag{13}
\]

If \(1<p<2\), then \(p'>2\), and applying \(12\) at exponent \(p'\) gives

\[
\boxed{
\mathfrak R_p
\le
2^{-1/p'}\sqrt{p'}.
}
\]

Thus the full open range \(1<p<\infty\), rather than only \(p>2\), is admissible.

\begin{enumerate}
\item Failure at \(p=\infty\) and \(p=1\)
\end{enumerate}
The open range is sharp.
For \(j=1,\dots,N\), set

\[
I_j=[0,2^{1-j}],
\qquad
I_{N+1}=[0,2^{-N}],
\]

and take

\[
f_j=\mathbf1_{I_j\setminus I_{j+1}}.
\]

The supports of the \(f_j\)'s are pairwise disjoint, so

\[
\left\|
\left(\sum_{j=1}^N|f_j|^2\right)^{1/2}
\right\|_\infty
=1.
\]

Moreover,

\[
|I_j\setminus I_{j+1}|=\frac12|I_j|,
\]

and hence

\[
P_{I_j}f_j=\frac12\mathbf1_{I_j}.
\]

On the smallest interval \(I_N\), every one of the \(N\) output functions is equal to \(1/2\). Thus

\[
\left(
\sum_{j=1}^N|P_{I_j}f_j(x)|^2
\right)^{1/2}
=
\frac{\sqrt N}{2},
\qquad x\in I_N.
\]

Therefore,

\[
\left\|
\left(\sum_{j=1}^N|P_{I_j}f_j|^2\right)^{1/2}
\right\|_\infty
=
\frac{\sqrt N}{2}.
\]

No constant independent of \(N\) can exist at \(p=\infty\).
If a uniform inequality held at \(p=1\), the corresponding operators \(T_{\mathcal I}\) would have uniformly bounded norms on \(L^1(\ell_2^N)\). Their adjoints, which are the same operators, would then be uniformly bounded on \(L^\infty(\ell_2^N)\), contradicting the preceding example. Hence the inequality also fails at \(p=1\).

\begin{enumerate}
\item Size of the optimal constants
\end{enumerate}
The preceding nested-interval example also shows that the upper bounds have the correct order as \(p\) approaches either endpoint.
On the annulus \(I_k\setminus I_{k+1}\), precisely the first \(k\) output functions are nonzero. Hence

\[
\left(
\sum_{j=1}^N|P_{I_j}f_j|^2
\right)^{1/2}
=
\frac{\sqrt k}{2}
\quad\text{on }I_k\setminus I_{k+1},
\]

whose measure is \(2^{-k}\). Therefore, for every \(N\ge k\),

\[
\mathfrak R_p
\ge
\frac12\sqrt k\,2^{-k/p}.
\tag{14}
\]

For \(p\ge2\), choose

\[
k=\left\lfloor\frac{p}{2\log 2}\right\rfloor.
\]

Since \(k\ge p/(4\log2)\) and \(k/p\le1/(2\log2)\),

\[
2^{-k/p}\ge e^{-1/2}.
\]

Thus

\[
\mathfrak R_p
\ge
\frac{e^{-1/2}}{4\sqrt{\log2}}\sqrt p.
\]

Using \(\mathfrak R_p=\mathfrak R_{p'}\) for \(p<2\), there is an absolute constant \(c>0\) such that

\[
c\sqrt{\max\{p,p'\}}
\le
\mathfrak R_p
\le
\sqrt{\max\{p,p'\}},
\qquad 1<p<\infty.
\]

Thus

\[
\mathfrak R_p\asymp
\sqrt{\max\{p,p'\}},
\]

up to absolute multiplicative constants. In particular, its divergence at \(p=1\) and \(p=\infty\) is quantitatively unavoidable.

Conclusion
The family of averaging projections onto arbitrary subintervals of \([0,1]\) satisfies the required square-function estimate for every \(p>2\), indeed for every \(1<p<\infty\). An explicit admissible choice in the range asked for in Q899 is

\[
\boxed{
C(p)=2^{-1/p}\sqrt p
\qquad (p>2).
}
\]

The complete exponent range is \(1<p<\infty\); the two endpoints fail.
